\begin{table}
\begin{tcolorbox}[tabularx={X|p{2cm}|X|},enhanced,fonttitle=\bfseries\large,fontupper=\normalsize\sffamily,
colback=yellow!10!white,colframe=red!50!black,colbacktitle=blue!30!white,
coltitle=black,title=Creating Varieties]
\hline
Command & Arguments & Purpose \\
\hline
\hline
\verb'-label_txt' & label & Specify an ASCII label for the variety. \\
\hline
\verb'-label_tex' & label & Specify a tex label for the variety. \\
\hline
\verb'-projective_space' & $P$ & Specify the underlying projective space. \\
\hline
\verb'-ring' & $R$ & Specify the underlying polynomial ring. \\
\hline
\verb'-equation_in_algebraic_form' & eqn & The defining equation in algebraic form. \\
\hline
\verb'-equation_by_coefficients' & eqn & The defining equation givem as a list of coefficients in the chosen ordering of monomials. \\
\hline
\verb'-second_equation_in_algebraic_form' & eqn & An optional secondary defining equation in algebraic form. \\
\hline
\verb'-second_equation_by_coefficients' & eqn & An optional secondary defining equation as a list of coefficients in the chosen ordering of monomials. \\
\hline
\verb'-points' & pts & The set of $\bbF_q$-rational points (optional). \\
\hline
\verb'-bitangents' & lines & A set of lines (optional). \\
\hline
\end{tcolorbox}
\caption{\label{tab:variety}Creating Varieties}
\end{table}
